Na camada de infraestrutura, o primeiro objetivo específico foi atingido: o nó sensor, o \textit{gateway} e o \textit{backend} foram construídos e operam em conjunto sobre um enlace LoRa ponto a ponto, e cada quadro recebido é registrado com os metadados de recepção, nos mesmos campos das variáveis candidatas dos modelos de detecção. Quatro coletas em bancada, descritas na Seção~\ref{sec:resultados_infra}, sustentam a validação. Na de referência, de 12 para 13 de outubro de 2026, o \textit{gateway} recebeu 6.007 quadros em 8,3 horas sem nenhuma perda de rádio e sem nenhuma rejeição, com intervalo médio entre quadros de 4.999,99~ms e desvio padrão de 0,80~ms para um período de 5~s. Nessa mesma coleta a rede entre \textit{gateway} e \textit{backend} caiu e não voltou, e os 3.575 quadros recebidos nas cinco horas seguintes foram recuperados do registro serial, sem lacuna nem duplicata: a coleta não depende da rede. Os dois caminhos de segurança, AES-128-GCM e Ascon-128a, ocupam o mesmo quadro de 26 bytes, com o mesmo tempo no ar de 61,7~ms, e a troca entre eles é feita na compilação, sem alterar o restante do firmware. Em operação, a decifragem de um quadro levou em média 215~$\mu$s no \textit{gateway}, 0,3\% do tempo no ar. Dos 23 requisitos da infraestrutura, 21 foram verificados e dois ficaram atendidos em parte (Subseção~\ref{subsec:rastreabilidade_infra}).

Os limites desses resultados são os da bancada. Os dados de aplicação são sintéticos e o nó não entra em modo de repouso entre transmissões, o que deixa parciais os requisitos RF1 e RNF1. As coletas registram apenas operação normal, com as placas a poucos metros uma da outra: o SNR fica saturado, o RSSI varia poucos decibéis e nenhum tráfego anômalo foi gerado. Elas servem, portanto, como referência de comportamento normal para a frente de detecção, e não como conjunto de treinamento. O consumo de energia do nó e o alcance do enlace não foram medidos.
